\documentclass[11pt,a4paper]{article}
\usepackage[utf8]{inputenc}
\usepackage{amsmath,amssymb,amsthm}
\usepackage{listings}
\usepackage{xcolor}
\usepackage[margin=2.5cm]{geometry}
\usepackage{hyperref}

\lstset{
  language=C++,
  basicstyle=\ttfamily\small,
  keywordstyle=\color{blue}\bfseries,
  commentstyle=\color{green!60!black},
  stringstyle=\color{red!70!black},
  numbers=left,
  numberstyle=\tiny\color{gray},
  breaklines=true,
  frame=single,
  tabsize=4,
  showstringspaces=false
}

\title{IOI 1999 -- Flower Shop}
\author{Solution}
\date{}

\begin{document}
\maketitle

\section{Problem Statement}

There are $F$ flowers and $V$ vases in a row ($F \le V$).
Each flower $i$ has a preference value $\mathrm{pref}[i][j]$ for vase $j$.
Flowers must be placed in \emph{strictly increasing} vase order:
$v_1 < v_2 < \cdots < v_F$.
Maximize the total preference and output the assignment.

\section{Solution Approach}

\subsection{DP Formulation}

Define:
\[
  \mathrm{dp}[i][j] = \max \text{ total preference placing flowers } 1, \ldots, i
                       \text{ with flower } i \text{ in vase } j.
\]

\textbf{Base case:} $\mathrm{dp}[1][j] = \mathrm{pref}[1][j]$ for
$1 \le j \le V - F + 1$.

\textbf{Transition:}
\[
  \mathrm{dp}[i][j] = \max_{k < j} \mathrm{dp}[i{-}1][k] + \mathrm{pref}[i][j],
  \quad i \le j \le V - F + i.
\]

\subsection{Running Maximum Optimization}

Rather than iterating over all $k < j$ for each $(i, j)$, we maintain a running maximum
of $\mathrm{dp}[i{-}1][k]$ as $j$ increases. This reduces the time per flower from
$O(V^2)$ to $O(V)$.

\subsection{Reconstruction}

Store the vase $k$ achieving the maximum in a parent array and trace backward from
the optimal final assignment.

\section{C++ Solution}

\begin{lstlisting}
#include <cstdio>
#include <cstring>
#include <algorithm>
#include <vector>
using namespace std;

int main() {
    int F, V;
    scanf("%d %d", &F, &V);

    vector<vector<int>> pref(F + 1, vector<int>(V + 1));
    for (int i = 1; i <= F; i++)
        for (int j = 1; j <= V; j++)
            scanf("%d", &pref[i][j]);

    const int NEG_INF = -1000000000;
    vector<vector<int>> dp(F + 1, vector<int>(V + 1, NEG_INF));
    vector<vector<int>> par(F + 1, vector<int>(V + 1, 0));

    // Base case: flower 1
    for (int j = 1; j <= V - F + 1; j++)
        dp[1][j] = pref[1][j];

    // Fill DP with running maximum
    for (int i = 2; i <= F; i++) {
        int best = NEG_INF;
        int bestk = 0;
        for (int j = i; j <= V - F + i; j++) {
            // Update running max from dp[i-1][j-1]
            if (dp[i - 1][j - 1] > best) {
                best = dp[i - 1][j - 1];
                bestk = j - 1;
            }
            if (best > NEG_INF) {
                dp[i][j] = best + pref[i][j];
                par[i][j] = bestk;
            }
        }
    }

    // Find the optimal vase for the last flower
    int ans = NEG_INF, lastVase = 0;
    for (int j = F; j <= V; j++) {
        if (dp[F][j] > ans) {
            ans = dp[F][j];
            lastVase = j;
        }
    }

    printf("%d\n", ans);

    // Reconstruct the assignment
    vector<int> assignment(F + 1);
    assignment[F] = lastVase;
    for (int i = F; i >= 2; i--)
        assignment[i - 1] = par[i][assignment[i]];

    for (int i = 1; i <= F; i++)
        printf("%d%c", assignment[i], i < F ? ' ' : '\n');

    return 0;
}
\end{lstlisting}

\section{Correctness}

The DP processes flowers in order $1, 2, \ldots, F$. For each flower $i$, the running
maximum ensures that $\mathrm{dp}[i][j]$ considers all valid predecessor vases $k < j$.
The constraint $i \le j \le V - F + i$ ensures that enough vases remain for subsequent
flowers.

\section{Complexity Analysis}

\begin{itemize}
  \item \textbf{Time complexity:} $O(F \times V)$. Each flower iterates over at most $V$
        vases, and the running maximum makes each iteration $O(1)$.
  \item \textbf{Space complexity:} $O(F \times V)$ for the DP and parent tables.
\end{itemize}

\end{document}
